% t2b_conjecture_note.tex — The Transcendental Ratio Identity
% A standalone conjecture note for Zenodo priority posting.
\documentclass[11pt]{amsart}

\usepackage[T1]{fontenc}
\usepackage{lmodern}
\usepackage{microtype}
\usepackage{amsmath,amssymb,amsthm}
\usepackage{mathtools}
\usepackage{booktabs}
\usepackage{enumitem}
\usepackage{hyperref}
\hypersetup{colorlinks=true, linkcolor=blue, citecolor=blue, urlcolor=blue}

\newtheorem{theorem}{Theorem}[section]
\newtheorem{proposition}[theorem]{Proposition}
\newtheorem{lemma}[theorem]{Lemma}
\newtheorem{conjecture}[theorem]{Conjecture}
\theoremstyle{definition}
\newtheorem{definition}[theorem]{Definition}
\theoremstyle{remark}
\newtheorem{remark}[theorem]{Remark}
\newtheorem{problem}{Problem}

\newcommand{\FF}{\mathcal{F}}
\newcommand{\KK}{\mathcal{K}}
\newcommand{\Rat}{\mathrm{Rat}}
\newcommand{\Log}{\mathrm{Log}}
\newcommand{\Alg}{\mathrm{Alg}}
\newcommand{\Trans}{\mathrm{Trans}}
\newcommand{\Des}{\mathrm{Des}}
\newcommand{\dps}{\mathrm{dps}}

\emergencystretch=1.5em

\title[The Transcendental Ratio Identity]{The Transcendental Ratio
Identity: A Conjecture on Degree-$(2,1)$ Polynomial Continued Fractions}

\author{Papanokechi}
\address{Independent researcher, Yokohama, Japan}
\email{shkubo@outlook.jp}

\subjclass[2020]{11A55, 11J70, 11J81, 11Y65, 40A15}

\keywords{Polynomial continued fractions, arithmetic stratification,
indicial equation, Worpitzky boundary, transcendental numbers,
experimental mathematics, conjecture}

\begin{document}

\begin{abstract}
For an integer polynomial continued fraction
$K = b_0 + \cfrac{a(1)}{b(1) + \cfrac{a(2)}{b(2) + \cdots}}$
of partial-numerator degree $2$ and partial-denominator degree $1$,
write $a(n) = a_2 n^2 + a_1 n + a_0$ and $b(n) = b_1 n + b_0$ with
$a_2 \ne 0$, $b_1 \ne 0$.  The asymptotic Worpitzky parameter
$r := a_2 / b_1^2$ stratifies $K$ via the indicial equation
$\lambda^2 - \lambda - r = 0$.  We propose the
\emph{Transcendental Ratio Identity}: every transcendental-stratum
limit lying in the negative-Worpitzky interior must arise at
$r = -2/9$ (indicial exponents $\{1/3, 2/3\}$).  The conjecture is
supported by zero counterexamples in $\sim 150{,}000$ families
across four independent experiments, including a complete census of
$\FF(2,5)$ and a targeted falsification at the three rationally
indicial $r \in \{-3/16, -4/25, -6/25\}$.  A separate Brouncker
boundary stratum at $r = +1/4$ is identified and rigorously
distinguished.  We outline a possible Schwarz-monodromic mechanism
and list four open problems for which the conjecture would be
decisive.
\end{abstract}

\maketitle

%% ===================================================================
\section{Introduction}
%% ===================================================================

A \emph{polynomial continued fraction (PCF) of degree $(d, e)$ with
integer coefficients} is an object
\begin{equation}\label{eq:pcf}
  K(a, b) \;=\; b_0 \,+\, \cfrac{a(1)}{b(1) \,+\, \cfrac{a(2)}{b(2)
  \,+\, \cfrac{a(3)}{b(3) \,+\, \cdots}}}
\end{equation}
where
\(
a(n) = a_d n^d + \cdots + a_0
\)
and
\(
b(n) = b_e n^e + \cdots + b_0
\)
are integer polynomials with $a_d \ne 0$ and $b_e \ne 0$.  The case
$(d,e) = (2,1)$ is the smallest non-trivial case in which classical
limits such as $\pi$, Brouncker's $4/\pi$, and the Ap\'ery-like series
all appear.

For $(d,e) = (2,1)$ we write $a(n) = a_2 n^2 + a_1 n + a_0$,
$b(n) = b_1 n + b_0$, and call
\begin{equation}\label{eq:r-defn}
  r \;:=\; \frac{a_2}{b_1^2}
\end{equation}
the \emph{Worpitzky parameter} of $K$.  The classical Worpitzky-Thron
theory of limit-periodic continued fractions identifies the disc
$r \in (-1/4,\,0)$ as the \emph{convergent interior} (negative
Worpitzky region), and the boundary points $r = -1/4$ and $r = +1/4$
as the parabolic boundary at which exotic phenomena (Brouncker's
$4/\pi$, the Lentz parabolic theorem) occur.

Empirical work begun in \cite{PCF25} surveying every integer PCF in
$\FF(2,4) := \{(d,e)=(2,1) : |a_i|, |b_j| \le 4\}$ uncovered an
unexpected algebraic identity:
\begin{quote}
\itshape Every one of the $24$ Trans-stratum limits in $\FF(2,4)$
satisfies $r = -2/9$.
\end{quote}
The corresponding indicial equation $\lambda^2 - \lambda - r = 0$
has roots $\{1/3,\,2/3\}$ at $r = -2/9$ --- a strikingly small
discrete value.  No analytic theory of which we are aware predicts
this constant.  Initial speculation that $-2/9$ might be a
``pigeonhole artifact'' (the unique rational of small height in the
negative-Worpitzky interior achievable at $D = 4$) motivated the
present work.

In this note we summarise four independent computational experiments
which together survey roughly $585{,}000$ degree-$(2,1)$ PCFs and
find \emph{zero} interior-stratum Trans counterexamples to the
relation $r = -2/9$.  We isolate a separate parabolic-boundary stratum at
$r = +1/4$ (the Brouncker class) which lies on the \emph{positive}
boundary and which the conjecture explicitly excludes.  We then
state the central conjecture.

\begin{conjecture}[Transcendental Ratio Identity]\label{conj:tri}
Let $K(a, b)$ be a degree-$(2,1)$ polynomial continued fraction with
integer coefficients, $a_2 \ne 0$, $b_1 \ne 0$.  If $K$ converges to
a limit $L$ in the Trans stratum --- that is, $L$ is not rational and
not expressible as a $\mathbb Q$-linear combination of logarithms of
algebraic numbers --- then
\begin{equation}\label{eq:conj}
  \frac{a_2}{b_1^2} \;=\; -\frac{2}{9}.
\end{equation}
Equivalently, the three-term recurrence associated with $K$ has
indicial exponents $\{1/3, 2/3\}$ at infinity.
\end{conjecture}

The remainder of the note proceeds as follows.
Section~\ref{sec:framework} gives the indicial-equation framework
and explains the Worpitzky interior/boundary distinction.
Section~\ref{sec:evidence} summarises the four computational
experiments supporting Conjecture~\ref{conj:tri}.
Section~\ref{sec:brouncker} treats the Brouncker boundary class
$r = +1/4$ and explains why it is not a counterexample.
Section~\ref{sec:open} lists candidate falsifiers, the open status
of the conjecture, and a speculative connection to Schwarz's list.

%% ===================================================================
\section{The Indicial Equation Framework}\label{sec:framework}
%% ===================================================================

The PCF~\eqref{eq:pcf} is governed by the three-term recurrence
\begin{equation}\label{eq:recurrence}
  y_{n+1} \;=\; b(n)\, y_n \,+\, a(n)\, y_{n-1},
  \qquad n \ge 1,
\end{equation}
applied separately to numerator partial-fraction sequences
$P_n$ and denominator sequences $Q_n$.  Setting $y_n = b_1^n n!\, z_n$
and analysing the leading-order behaviour, one obtains the indicial
equation
\begin{equation}\label{eq:indicial}
  \lambda^2 \,-\, \lambda \,-\, r \;=\; 0,
  \qquad r = \frac{a_2}{b_1^2},
\end{equation}
whose roots $\lambda_\pm = \tfrac{1}{2}\bigl(1 \pm \sqrt{1 + 4r}\bigr)$
encode the asymptotic exponents of the two formal solutions of
\eqref{eq:recurrence} at infinity.

\begin{lemma}\label{lem:strata}
The discriminant $\Delta = 1 + 4r$ partitions the parameter line
into the following regimes.
\begin{enumerate}[label={\upshape(\roman*)}]
\item $\Delta < 0$, i.e.\ $r < -1/4$: complex-conjugate exponents,
limit-periodic divergence in general.
\item $\Delta = 0$, i.e.\ $r = -1/4$: the negative parabolic
boundary; double root $\lambda = 1/2$.
\item $0 < \Delta < 1$, i.e.\ $-1/4 < r < 0$: the
\emph{negative-Worpitzky interior}; real exponents
$\lambda_\pm \in (0,1)$.
\item $\Delta = 1$, i.e.\ $r = 0$: the trivial linear case
$a(n)/b(n) \to 0$.
\item $\Delta > 1$, i.e.\ $r > 0$: positive Worpitzky region;
$\lambda_+ > 1$, $\lambda_- < 0$.  The boundary point $r = +1/4$
($\Delta = 2$) is the classical Brouncker locus.
\end{enumerate}
\end{lemma}

For $r = -2/9$ one has $\Delta = 1/9$ and indicial roots
$\{1/3,\,2/3\}$.  These are the simplest pair of \emph{rational}
exponents in the open interval $(0,1)$: writing $\lambda_\pm = p/q$
with $\gcd(p,q) = 1$ and $0 < p < q$, the ordered list of admissible
$(\lambda_-, \lambda_+) = (p/q,\,(q-p)/q)$ pairs in the negative
Worpitzky interior begins
\begin{align*}
  \{\tfrac{1}{3},\tfrac{2}{3}\} &\;\to\; r = -\tfrac{2}{9}, &
  \{\tfrac{1}{4},\tfrac{3}{4}\} &\;\to\; r = -\tfrac{3}{16}, \\
  \{\tfrac{1}{5},\tfrac{4}{5}\} &\;\to\; r = -\tfrac{4}{25}, &
  \{\tfrac{2}{5},\tfrac{3}{5}\} &\;\to\; r = -\tfrac{6}{25},
\end{align*}
followed by shifted variants at $q = 6$ and the prime-denominator
family at $q = 7$.  Conjecture~\ref{conj:tri} asserts that --- among
all of these admissible pairs --- only the first one supports a
Trans-stratum limit.

\begin{remark}[speculative]
\label{rem:schwarz-speculation}
The triangle group $\Gamma(3,3,\infty)$ uniformising the indicial
roots $\{1/3, 2/3\}$ is the smallest hyperbolic Hecke triangle group
at which an integer arithmetic structure exists.  We speculate that
Conjecture~\ref{conj:tri} reflects a Schwarz-monodromic constraint
on the second-order linear recurrence~\eqref{eq:recurrence}: the
formal monodromy at infinity must lift to an integer-coefficient
operator, and only $\{1/3,2/3\}$ among rational pairs in $(0,1) \times
(0,1)$ produces a fundamental domain of small enough discriminant for
this lift to be obstruction-free.  This is conjectural; we offer no
proof.  See Section~\ref{sec:open} for a precise open problem.
\end{remark}

%% ===================================================================
\section{Empirical Evidence}\label{sec:evidence}
%% ===================================================================

We summarise four computational experiments.  Each was logged as a
session in the SIARC relay bridge (commits referenced in the
disclosure section); raw results are reproducible from the published
scripts.

\subsection*{Experiment A: $\FF(2,4)$ complete census}
The base-case paper~\cite{PCF25} enumerated every integer
degree-$(2,1)$ PCF with $|a_i|, |b_j| \le 4$, totaling $531{,}441$
families, and identified $24$ Trans-stratum limits, all with
$r = -2/9$.

\subsection*{Experiment B: $\FF(2,5)$ complete census}
Extending the search to $D = 5$ ($133{,}100$ candidate families), the
F$(2,5)$ census~\cite{F25} found $88{,}224$ convergent families with
the stratum decomposition $\Rat(975) \sqcup \Alg(9{,}667) \sqcup
\Trans(70) \sqcup \Log(12) \sqcup \Des(77{,}500)$.  The Trans and
Log strata partition by Worpitzky parameter as follows:
\begin{itemize}
\item $r = -2/9$: $56$ Trans (M\"obius transforms of $\pi$) and
all $12$ Log (M\"obius transforms of $\log 2$) families;
\item $r = +1/4$: $14$ Trans families of Brouncker type;
\item every other Worpitzky parameter: $0$ Trans, $0$ Log.
\end{itemize}
Crucially, all $12$ Log-stratum families also exhibit $r = -2/9$:
the identity is independent of which transcendental constant
($\pi$ or $\log 2$) the limit involves.  This refutes any pigeonhole
explanation tied to a particular transcendental.

\subsection*{Experiment C: targeted resonance search}
The four ratios $r \in \{-3/16,\,-4/25,\,-6/25,\,-2/49\}$ are the
next admissible rational-indicial Worpitzky-interior values
(Section~\ref{sec:framework}).  Restricted enumeration at
$|b_1| \in \{4, 5\}$ (covering the first three ratios), with the
remaining slots $|a_1|, |a_0|, |b_0| \le 5$, was processed through
the same PSLQ-based three-stage pipeline used in Experiments A and
B.  Two follow-up resonance sessions extended the sweep to
$|b_1| \in \{6, 7\}$ and $|b_1| \in \{8,9,10,11,12\}$
(the latter covering the rational-indicial targets $r = -2/49$,
$-2/9$ at $b_1 = \pm 9, \pm 12$, and $r = -35/144$ at
$b_1 = \pm 12$).  Results are summarised in
Table~\ref{tab:resonance}.

\subsection*{Experiment D: indicial discriminant pre-screen}
A theoretical analysis of all rationals $r$ with $\sqrt{1 + 4r} \in
\mathbb Q$ at $|b_1| \le 5$ enumerates exactly ten admissible interior
ratios.  Of the ten, only $r = -2/9$ is realised in the Trans stratum
across Experiments A--C; the remaining nine are confirmed empty in
either A, B, or C.

\begin{table}[t]
\centering
\caption{Stratum counts in $\FF(2,4)$ and $\FF(2,5)$.}
\label{tab:censuses}
\begin{tabular}{lrrl}
\toprule
Search            & Trans & Log & Worpitzky parameter $r$ \\
\midrule
$\FF(2,4)$ ($D=4$)         & $24$ & $0$  & $-2/9$         \\
$\FF(2,5)$ ($D=5$, neg.\ interior) & $56$ & $12$ & $-2/9$         \\
$\FF(2,5)$ ($D=5$, Brouncker)      & $14$ & $0$  & $+1/4$         \\
\bottomrule
\end{tabular}
\end{table}

\begin{table}[t]
\centering
\caption{Targeted resonance searches.  Top block: $b_1 \in \{\pm 4,
\pm 5\}$, $|a_1|,|a_0|,|b_0| \le 5$.  Bottom block: extended
sweeps at higher $|b_1|$.  All runs at dps $=100$.}
\label{tab:resonance}
\begin{tabular}{lrrr}
\toprule
Scope (target ratio $r$ or range) & Convergent families & Trans & Log \\
\midrule
$b_1 \in \{\pm 4, \pm 5\}$, $r = -3/16$ & $2{,}384$ & $0$ & $0$ \\
$b_1 \in \{\pm 4, \pm 5\}$, $r = -4/25$ & $2{,}396$ & $0$ & $0$ \\
$b_1 \in \{\pm 4, \pm 5\}$, $r = -6/25$ & $2{,}394$ & $0$ & $0$ \\
\midrule
$b_1 \in \{\pm 6, \pm 7\}$ (full sweep) & $175{,}686$ & $0$ & $2$ \\
$b_1 \in \{\pm 8,\dots,\pm 12\}$ (motivated targets) & $259{,}280$ & $0$ & $0$ \\
\midrule
Combined        & $442{,}140$ & $0$ & $2$ \\
\bottomrule
\end{tabular}
\end{table}

\subsection*{Aggregate}
The combined empirical base across the experiments comprises
approximately $585{,}000$ integer degree-$(2,1)$ PCF families
surveyed (convergent count, including the extended
$b_1 \in \{6,7\}$ and $b_1 \in \{8,\dots,12\}$ resonance sweeps).
The two Log-stratum families found at $b_1 = \pm 6$ both lie at
$r = -1/36$ (a Bauer--Stern doubling parameter outside the
admissible rational-indicial set predicted by the conjecture's
framework) and are tracked under Section~\ref{sec:open}.  Notably,
at $b_1 = \pm 12$ the ratio $r = -35/144$ was tested explicitly
and produced $0$ hits.  The number of Trans-stratum families with
$r \ne -2/9$ \emph{and} $r$ in the negative Worpitzky
interior
(equivalently: the number of counterexamples to
Conjecture~\ref{conj:tri}) is exactly zero.  The $14$ Brouncker
families at $r = +1/4$ lie on the positive Worpitzky boundary and
are addressed in Section~\ref{sec:brouncker}.

%% ===================================================================
\section{The Brouncker Boundary Class}\label{sec:brouncker}
%% ===================================================================

The $14$ families at $r = +1/4$ that appear in $\FF(2,5)$ are
\emph{not} counterexamples to Conjecture~\ref{conj:tri}.  We
distinguish them as follows.

\begin{proposition}\label{prop:brouncker-boundary}
The Worpitzky parameter $r = +1/4$ lies on the positive parabolic
boundary $\Delta = 2$, with indicial double root $\lambda = 1/2$
realised in $\mathbb{R}$ but lying \emph{outside} the negative
Worpitzky interior $r \in (-1/4, 0)$ to which
Conjecture~\ref{conj:tri} applies.
\end{proposition}

A representative Brouncker family $(a_2, a_1, a_0, b_1, b_0) =
(1, 0, 0, -2, -1)$ produces, at integer recurrence, the limit
$L = -4/\pi$, recovering Lord Brouncker's celebrated
$1655$ identity~\cite{Brouncker} after sign normalisation.  The full
$14$-element class consists of small integer perturbations of
Brouncker's recurrence; PSLQ identifies all $14$ as M\"obius
transforms of $4/\pi$ at residual below $10^{-30}$.

The boundary nature of $r = +1/4$ is essential: the Brouncker
asymptotics are parabolic (algebraic decay of partial-fraction error),
whereas Conjecture~\ref{conj:tri}'s domain --- the strictly negative
Worpitzky interior --- is hyperbolic (geometric decay).  These are
qualitatively different convergence regimes.

\begin{remark}
The negative parabolic boundary $r = -1/4$ has \emph{no integer
realisation} at $|b_1| \le 5$ (the equation $a_2 = -b_1^2/4$ has no
integer solutions for $b_1 \in \{1,\dots,5\}$); the question of
whether a parabolic-negative analogue of the Brouncker class exists
at larger $|b_1|$ is open.
\end{remark}

%% ===================================================================
\section{Candidate Falsifiers and Open Problems}\label{sec:open}
%% ===================================================================

Four candidate falsifiers and four open problems.

\subsection*{Candidate falsifiers (status)}
\begin{enumerate}[label={\upshape(\arabic*)}]
\item $r = -3/16$, indicial roots $\{1/4, 3/4\}$, requires $|b_1| = 4$.
\textbf{Empty in Trans/Log at $|b_1| = 4$, $|a_i|, |b_j| \le 5$.}
\item $r = -4/25$, indicial roots $\{1/5, 4/5\}$, requires
$|b_1| = 5$. \textbf{Empty.}
\item $r = -6/25$, indicial roots $\{2/5, 3/5\}$, requires
$|b_1| = 5$. \textbf{Empty.}
\item $r = -2/49$, indicial roots $\{1/7, 6/7\}$, requires
$|b_1| = 7$.  \textbf{Untested at the time of writing} (outside the
$|b_1| \le 5$ scope of Experiments A--C); the natural next
falsification target.
\end{enumerate}

\subsection*{Open Problems}
\begin{problem}\label{prob:d6}
Does Conjecture~\ref{conj:tri} survive at $D = 6$ and $D = 7$?
The natural extension is a complete $\FF(2,7)$ census; the search
size grows as $(2D{+}1)^5$, well within tractable PSLQ budgets at
$\dps = 100$.
\end{problem}

\begin{problem}\label{prob:rational-coeffs}
Does the analogue of Conjecture~\ref{conj:tri} hold for
\emph{rational}-coefficient PCFs of degree $(2,1)$, with $r$ defined
as in~\eqref{eq:r-defn}?  Equivalently, is the integer-coefficient
restriction essential?
\end{problem}

\begin{problem}\label{prob:schwarz}
Is there a Schwarz-monodromic proof of Conjecture~\ref{conj:tri}
along the lines sketched in Remark~\ref{rem:schwarz-speculation}?
A starting point would be a classification of integer-coefficient
second-order linear recurrences with rational indicial exponents
$\{p/q,\,(q-p)/q\}$, $q \le 7$, having a transcendental ratio
$P_n/Q_n$ as $n \to \infty$.
\end{problem}

\begin{problem}\label{prob:negative-parabolic}
Does a negative parabolic boundary stratum exist at $r = -1/4$?
The integer constraint $a_2 = -b_1^2/4$ forces $b_1$ even and
$|b_1| \ge 2\sqrt{|a_2|}$; the smallest candidate is
$(a_2, b_1) = (-1, 2)$.  An empirical survey is straightforward.
\end{problem}

\subsection*{Closing remark}
Conjecture~\ref{conj:tri} is, to the best of our knowledge, the
first published quantitative arithmetic constraint on the
transcendental stratum of degree-$(2,1)$ integer PCFs that is
\emph{independent of the transcendental constant} appearing in the
limit.  Its proof or refutation, in either the integer or rational
setting, would constitute a substantive advance in the
arithmetic theory of polynomial continued fractions.

%% ===================================================================
\section*{Computational and AI-Assistance Disclosure}
%% ===================================================================

\textbf{Reproducibility.}  All computations underlying this note are
reproducible from publicly archived scripts and datasets.  The four
experiments referenced in Section~\ref{sec:evidence} correspond to
sessions in the SIARC relay bridge repository
\href{https://github.com/papanokechi/siarc-relay-bridge}%
{papanokechi/siarc-relay-bridge}, in particular sessions
\texttt{2026-04-29/T2B-APERY-INVESTIGATION} (commit
\texttt{0161a321}; theoretical pre-analysis),
\texttt{2026-04-29/T2B-F25-FALSIFICATION} (commit
\texttt{6e28269}; the F$(2,5)$ census), and
\texttt{2026-04-29/T2B-RESONANCE-SEARCH} (commit
\texttt{ab0d70b}; the targeted resonance test).  Each session is
accompanied by a handoff document, AEAL claim entries, and bit-exact
numerical inputs/outputs.

\textbf{AEAL.}  Numerical claims in this note follow the
``Aggressive Empirical Audit Layer'' (AEAL) discipline: every quoted
count or ratio is logged as a \texttt{claims.jsonl} entry with
\texttt{evidence\_type = computation}, the
governing precision (\texttt{dps=100} for Experiments B and C;
\texttt{dps=300}--$500$ for Experiment A), the script filename,
and a SHA-256 hash of the result file.  See the bridge sessions for
the full ledger.

\textbf{AI-assistance disclosure.}  Drafting and verification of this
manuscript made use of two AI systems in supporting roles.
GitHub Copilot (Anthropic Claude backend) was used as the execution
agent for all enumerations, PSLQ classifications, and figure-data
generation; the agent's outputs are deterministic functions of the
stated scripts.  Anthropic Claude (\texttt{claude.ai}) was used as a
critique and drafting partner for the manuscript text.  All
mathematical claims in this note were verified by the author against
the underlying numerical data.

%% ===================================================================
\begin{thebibliography}{99}
%% ===================================================================

\bibitem{PCF25}
Papanokechi.
\emph{A Complete Arithmetic Stratification of Degree-2 Polynomial
Continued Fractions: The $\FF(2,4)$ Base Case.}
Submitted to Mathematics of Computation, 2026.
Preprint available at
\href{https://github.com/papanokechi/pcf-research}%
{github.com/papanokechi/pcf-research}.

\bibitem{F25}
Papanokechi.
\emph{F$(2,5)$ Census and the Pigeonhole Falsification.}
SIARC relay bridge session
\texttt{2026-04-29/T2B-F25-FALSIFICATION}, commit \texttt{6e28269},
2026.

\bibitem{ThronWaadeland}
W.~B.~Jones and W.~J.~Thron,
\emph{Continued Fractions: Analytic Theory and Applications,}
Encyclopedia of Mathematics and its Applications, vol.~11,
Addison-Wesley, Reading, MA, 1980.

\bibitem{Worpitzky}
J.~Worpitzky,
\emph{Untersuchungen \"uber die Entwickelung der monodromen und
monogenen Functionen durch Kettenbr\"uche,}
Programm Friedrichs-Gymnasium und Realschule, Berlin, 1865.

\bibitem{Brouncker}
W.~Brouncker, in J.~Wallis,
\emph{Arithmetica Infinitorum}, Oxford, 1656.
(Brouncker's continued fraction $\tfrac{4}{\pi}
= 1 + \cfrac{1^2}{2 + \cfrac{3^2}{2 + \cfrac{5^2}{2+\cdots}}}$.)

\bibitem{Schwarz}
H.~A.~Schwarz,
\emph{\"Uber diejenigen F\"alle, in welchen die Gauss\-ische
hypergeometrische Reihe eine algebraische Function ihres vierten
Elementes darstellt,}
J.\ Reine Angew.\ Math.\ \textbf{75} (1873), 292--335.

\bibitem{RamanujanMachine}
G.~Raayoni, S.~Gottlieb, Y.~Manor, G.~Pisha, Y.~Harris, U.~Mendlovic,
D.~Haviv, Y.~Hadad, and I.~Kaminer,
\emph{Generating conjectures on fundamental constants with the
Ramanujan Machine,}
Nature \textbf{590} (2021), 67--73.

\end{thebibliography}

\end{document}
